%% =========================================================
%% Chapter: Periodic MPS — applications of the irreducible-form
%% Fundamental Theorem.  Mirrors §4 of arXiv:1708.00029
%% (de las Cuevas--Schuch--Pérez-García--Cirac).
%% =========================================================

\chapter{Periodic MPS: physical symmetries and refinement}\label{ch:periodic_ft_apps}

This chapter records the two applications of the periodic equal-case
Fundamental Theorem of MPS (Theorem~\ref{thm:periodic_ft}) developed in
\cite{DeLasCuevas2017Irreducible}: the symmetry corollary lifting a
physical on-site symmetry of a tensor in irreducible form to a virtual
$Z$-gauge, and the equivalence between $p$-refinability of the
matrix-product-vector family and $p$-divisibility of the transfer map.

\section{Periodic irreducible form and physical-index rotation}%
\label{sec:periodic_irreducible_rotation}

\begin{definition}[Periodic tensor]
    \label{def:periodic_tensor}
    \lean{MPSTensor.IsPeriodic}
    \leanok
    An irreducible MPS tensor $A$ of bond dimension $D$ is
    \emph{$m$-periodic} (for some $m \ge 1$) if the peripheral spectrum
    of its transfer map $\E_A$ is exactly the $m$-th roots of unity
    $\{e^{2\pi i k/m} : 0 \le k < m\}$.
    A $1$-periodic tensor is exactly an irreducible tensor with
    primitive transfer map (Definition~\ref{def:primitive_channel}).
\end{definition}

\begin{definition}[Irreducible form]
    \label{def:irreducible_form}
    \lean{MPSTensor.IsIrreducibleForm}
    \leanok
    A tensor $A$ is in \emph{irreducible form} if it is block-diagonal
    \[
        A^i = \bigoplus_{j \in J} \mu_j\, A_j^i,
    \]
    where each block $A_j$ is an $m_j$-periodic tensor,
    each weight $\mu_j \neq 0$, and the blocks are pairwise
    non-repeated (no two distinct $j, j'$ have $A_j$ gauge-equivalent
    to $A_{j'}$).
    This is the standard form introduced in \cite[\S II]{DeLasCuevas2017Irreducible}.
    It extends the normal canonical form
    (Definition~\ref{def:is_normal_canonical_form}),
    which requires the stronger condition that every block is $1$-periodic.
\end{definition}

\begin{definition}[Physical-index rotation]%
    \label{def:rotate_physical}
    \lean{MPSTensor.rotatePhysical}
    \leanok
    \uses{def:mps_tensor}
    Given a matrix $u \in M_d(\C)$ and an MPS tensor~$A$,
    the \emph{physical-index rotation} is the tensor
    $(A_u)^i := \sum_j u_{ij}\, A^j$.
\end{definition}

\begin{theorem}[Transfer map preserved by physical rotation]%
    \label{thm:transfer_map_rotate_physical}
    \lean{MPSTensor.transferMap_rotatePhysical}
    \leanok
    \uses{def:rotate_physical, def:transfer_map}
    If $u u^\dagger = I$, then the transfer map is invariant under
    physical-index rotation:
    $\mathcal{E}_{A_u} = \mathcal{E}_A$.
\end{theorem}

\begin{proof}\leanok
    \uses{lem:unitary_kraus_mixing}
    Unitary freedom of the Kraus representation (Theorem~2.18,
    Wolf \emph{Quantum Channels \& Operations}).
\end{proof}

\begin{theorem}[Left-canonical preserved by physical rotation]%
    \label{thm:left_canonical_rotate_physical}
    \lean{MPSTensor.isLeftCanonical_rotatePhysical}
    \leanok
    \uses{def:rotate_physical}
    If $u u^\dagger = I$ and $A$ is left-canonical, then
    $A_u$ is left-canonical.
\end{theorem}

\begin{proof}\leanok
    \uses{def:rotate_physical}
    Expand $\sum_i B_i^\dagger B_i$ where $B_i = \sum_j u_{ij} A_j$,
    swap sums, apply orthogonality $u^\dagger u = I$, and collapse to
    $\sum_j A_j^\dagger A_j = I$.
\end{proof}

\begin{theorem}[Irreducibility preserved by physical rotation]%
    \label{thm:irreducible_tensor_rotate_physical}
    \lean{MPSTensor.isIrreducibleTensor_rotatePhysical}
    \leanok
    \uses{def:rotate_physical}
    If $u u^\dagger = I$ and $A$ is irreducible, then
    $A_u$ is irreducible.
\end{theorem}

\begin{proof}\leanok
    \uses{def:rotate_physical}
    Contrapositive: any nontrivial invariant projection~$P$ for the
    rotated tensor yields $(1-P)\, A_k\, P = 0$ for all~$k$ via the
    orthogonality of~$u$, contradicting irreducibility of~$A$.
\end{proof}

\begin{theorem}[Periodicity preserved by physical rotation]%
    \label{thm:periodic_rotate_physical}
    \lean{MPSTensor.isPeriodic_rotatePhysical}
    \leanok
    \uses{def:rotate_physical, thm:transfer_map_rotate_physical,
        thm:left_canonical_rotate_physical,
        thm:irreducible_tensor_rotate_physical}
    If $u u^\dagger = I$ and $A$ is periodic with period~$m$, then
    $A_u$ is periodic with the same period.
\end{theorem}

\begin{proof}\leanok
    \uses{thm:transfer_map_rotate_physical,
        thm:left_canonical_rotate_physical,
        thm:irreducible_tensor_rotate_physical}
    Assembles transfer-map invariance, left-canonical preservation, and
    irreducibility preservation.
\end{proof}

\begin{theorem}[SameMPV preserved by physical rotation]%
    \label{thm:same_mpv_rotate_physical}
    \lean{MPSTensor.sameMPV₂_rotatePhysical}
    \leanok
    \uses{def:rotate_physical, def:same_mpv}
    If $\mathcal{V}(A) = \mathcal{V}(B)$, then
    $\mathcal{V}(A_u) = \mathcal{V}(B_u)$.
\end{theorem}

\begin{proof}\leanok
    \uses{def:rotate_physical, def:same_mpv}
    Induction on the word length with a strengthened hypothesis
    carrying an arbitrary prefix matrix, reducing each step to the
    hypothesis via word concatenation.
\end{proof}

\begin{theorem}[Physical rotation distributes over block assembly]%
    \label{thm:rotate_physical_blocks}
    \lean{MPSTensor.rotatePhysical_toTensorFromBlocks}
    \leanok
    \uses{def:rotate_physical}
    The rotated block-diagonal tensor satisfies
    $(\bigoplus_k \mu_k A_k)_u
        = \bigoplus_k \mu_k (A_k)_u$.
\end{theorem}

\begin{proof}\leanok
    \uses{def:rotate_physical}
    Pointwise matrix-entry computation using linearity of
    the block-diagonal embedding and reindexing.
\end{proof}

\begin{theorem}[Irreducible form preserved by physical rotation]%
    \label{thm:irreducible_form_rotate_physical}
    \lean{MPSTensor.isIrreducibleForm_rotatePhysical}
    \leanok
    \uses{def:rotate_physical, def:irreducible_form,
        thm:periodic_rotate_physical,
        thm:same_mpv_rotate_physical,
        thm:rotate_physical_blocks}
    If $u u^\dagger = I$ and $A$ is in irreducible form~II, then
    $A_u$ is in irreducible form~II with the same block structure and
    rotated blocks.
\end{theorem}

\begin{proof}\leanok
    \uses{thm:periodic_rotate_physical,
        thm:same_mpv_rotate_physical,
        thm:rotate_physical_blocks}
    Each block remains periodic by
    Theorem~\ref{thm:periodic_rotate_physical}, the SameMPV condition
    transfers via Theorem~\ref{thm:same_mpv_rotate_physical}, and
    block-assembly compatibility follows from
    Theorem~\ref{thm:rotate_physical_blocks}.
\end{proof}

\subsection{The \texorpdfstring{$Z$}{Z}-gauge degree of freedom}

\begin{definition}[Entrywise periodic \(Z\)-gauge ratio]\label{def:z_gauge_entry}
    \lean{MPSTensor.zGaugeEntry}
    \leanok
    Given complex weights $\mu,\nu$, define the entrywise
    \(Z\)-gauge ratio by
    \[
        z(\mu,\nu) := \mu/\nu.
    \]
\end{definition}

\begin{lemma}[Matched powers imply root-of-unity ratio]\label{lem:z_gauge_entry_pow}
    \lean{MPSTensor.zGaugeEntry_pow_eq_one_of_pow_eq}
    \leanok
    \uses{def:z_gauge_entry}
    If $\mu^m=\nu^m$ and $\nu\neq 0$, then
    \[
        z(\mu,\nu)^m = 1.
    \]
\end{lemma}

\begin{lemma}[Ratio rescales denominator back to numerator]\label{lem:z_gauge_entry_mul}
    \lean{MPSTensor.zGaugeEntry_mul_right}
    \leanok
    \uses{def:z_gauge_entry}
    If $\nu\neq 0$, then
    \[
        z(\mu,\nu)\,\nu=\mu.
    \]
\end{lemma}

\begin{definition}[Diagonal periodic \(Z\)-gauge matrix]\label{def:z_gauge_diagonal}
    \lean{MPSTensor.zGaugeDiagonal}
    \leanok
    \uses{def:z_gauge_entry}
    For matched weight families $(\mu_i)_i$ and $(\nu_i)_i$ on a finite
    index type, define the diagonal matrix
    \[
        Z := \diag\!\bigl(z(\mu_i,\nu_i)\bigr)_i.
    \]
\end{definition}

\begin{lemma}[Diagonal \(Z\)-gauge has finite order under matched powers]\label{lem:z_gauge_diagonal_pow}
    \lean{MPSTensor.zGaugeDiagonal_pow_eq_one}
    \leanok
    \uses{def:z_gauge_diagonal, lem:z_gauge_entry_pow}
    If $\mu_i^m=\nu_i^m$ and $\nu_i\neq 0$ for all~$i$, then
    \[
        Z^m=\Id.
    \]
\end{lemma}

\begin{lemma}[Diagonal \(Z\)-gauge transports diagonal weights]\label{lem:z_gauge_diagonal_mul}
    \lean{MPSTensor.zGaugeDiagonal_mul_diagonal}
    \leanok
    \uses{def:z_gauge_diagonal, lem:z_gauge_entry_mul}
    Under $\nu_i\neq 0$ for all~$i$,
    \[
        Z\,\diag(\nu)=\diag(\mu).
    \]
\end{lemma}

\begin{definition}[$Z$-gauge equivalence]
    \label{def:z_gauge_equiv}
    \lean{MPSTensor.ZGaugeEquiv}
    \leanok
    Let $A$ be an $m$-periodic irreducible tensor of bond dimension $D$.
    A \emph{$Z$-gauge transformation} of $A$ is a diagonal unitary
    matrix $Z \in \MN{D}$ satisfying:
    \begin{enumerate}
        \item $Z$ commutes with every Kraus operator: $[A^i, Z] = 0$
              for all physical indices $i$, and
        \item $Z^m = \Id_D$ (all diagonal entries are $m$-th roots of unity).
    \end{enumerate}
    Replacing $A$ by $ZA$ (i.e.\ $A^i \mapsto Z A^i$) leaves the
    MPV family invariant for lengths $N$ that are multiples of $m$,
    because the factor $\tr(Z^N) = \tr(\Id) = D$ for $N \equiv 0 \pmod{m}$.
    Two tensors $A, B$ are \emph{$Z$-gauge equivalent} if there exists
    an invertible $Y$ and a $Z$-gauge transformation $Z$ for $A$ such that
    $Z A^i = Y B^i Y^{-1}$ for all~$i$.
\end{definition}

\begin{remark}[$Z$-gauge is trivial for period-$1$ blocks]
    \label{rem:z_gauge_trivial_period1}
    \lean{MPSTensor.ZGaugeEquiv.of_period_one}
    \leanok
    When $m = 1$ the condition $Z^1 = \Id$ forces $Z = \Id$,
    so $Z$-gauge equivalence reduces to ordinary gauge equivalence
    (Definition~\ref{def:gauge_equiv}).
    The $Z$-gauge is therefore a genuine new degree of freedom
    that appears only for $m > 1$.
\end{remark}

\section{Physical symmetries on periodic MPS}\label{sec:periodic_symmetry}

\begin{theorem}[Cor.~4.1 — Physical symmetry yields a virtual $Z$-gauge]%
    \label{cor:cor_4_1_physical_symmetry_zgauge}
    \lean{MPSTensor.cor_4_1_physical_symmetry_zgauge}
    \leanok
    \uses{def:irreducible_form, def:z_gauge_equiv, def:on_site_symmetric,
          def:twisted_tensor, thm:periodic_ft}
    Let $A$ be an MPS tensor in irreducible form~II and let
    $U : G \to \GL_d(\C)$ be a representation of a group $G$ on the
    physical leg, acting unitarily ($U(g)\, U(g)^{\dagger} = \Id$).
    Assume moreover the periodic equal-case Fundamental Theorem of MPS
    (Theorem~\ref{thm:periodic_ft}) as an explicit hypothesis on the
    pair $(d, D)$ — formalised in Lean as the proposition
    \lean{MPSTensor.PeriodicEqualCaseFT}\,$d\,D$.
    If $A$ is on-site symmetric under~$U$ — that is, the twisted tensor
    $\widetilde{A}_g$ has the same MPV family as~$A$ for every $g \in G$
    — then for each $g \in G$ there exists a positive period $m_g$ and
    a $\Z_{m_g}$-gauge equivalence between~$A$ and~$\widetilde{A}_g$.

    Concretely there are matrices $Z_g$ and~$Y_g$, with $Z_g^{m_g} = \Id$
    and $[A^i, Z_g] = 0$, satisfying
    \[
        Z_g\, A^i = Y_g\, \widetilde{A}_g^{\,i}\, Y_g^{-1}
        \quad\text{for every physical index } i.
    \]
\end{theorem}

\begin{proof}\leanok
    \uses{thm:periodic_ft, thm:irreducible_form_rotate_physical,
          cor:symmetry_periodic_assembly}
    The twisted tensor $\widetilde{A}_g$ coincides with the physical-index
    rotation $A_{U(g)}$ of Definition~\ref{def:rotate_physical}.
    By Theorem~\ref{thm:irreducible_form_rotate_physical} the rotated
    tensor remains in irreducible form~II, and the on-site symmetry
    hypothesis supplies the equality $\mathcal{V}(A) = \mathcal{V}(A_{U(g)})$.
    Apply the periodic equal-case Fundamental Theorem
    (Theorem~\ref{thm:periodic_ft}) to obtain the asserted
    $\Z_{m_g}$-gauge equivalence.
\end{proof}

\begin{definition}[Periodic projective-rigidity hypothesis]%
    \label{def:periodic_projective_rigidity}
    \lean{MPSTensor.PeriodicProjectiveRigidity}
    \leanok
    \uses{def:twisted_tensor, def:on_site_symmetric}
    Let $A : \mathrm{MPSTensor}\ d\ D$ be on-site symmetric under a group
    action $U : G \to \GL_d(\C)$.  The tensor $A$ satisfies the
    \emph{periodic projective-rigidity hypothesis} if one may choose a
    family of virtual gauges $Y : G \to \GL_D(\C)$ and a scalar
    $2$-cochain $u : G \times G \to \C^{\times}$ such that

    \begin{enumerate}
        \item every $Y_g$ realises the Corollary~4.1 intertwining with
              the twisted tensor $\widetilde{A}_g$ for some accompanying
              $Z_g$ (with $Z_g^{m_g} = \Id$ and $[A^i, Z_g] = 0$); and
        \item the family obeys the projective multiplication law on the
              bond space: $Y_h\, Y_g = u(g, h)\, Y_{gh}$.
    \end{enumerate}

    The injective case reduces the rigidity to scalar uniqueness of
    gauges (cf.~\S\ref{sec:injective_symmetry_cocycle}); in the
    genuinely periodic setting the bond-space commutant of $A$ is
    generated by the $Z_g$ matrices, so coherence of the factor system
    $u$ is an independent analytic input.
\end{definition}

\begin{theorem}[Cor.~4.1 projective-representation upgrade]%
    \label{thm:cor_4_1_projective_rep}
    \lean{MPSTensor.cor_4_1_projective_rep}
    \leanok
    \uses{cor:cor_4_1_physical_symmetry_zgauge,
          def:periodic_projective_rigidity}
    Under the hypotheses of Corollary~\ref{cor:cor_4_1_physical_symmetry_zgauge}
    together with the periodic projective-rigidity hypothesis
    (Definition~\ref{def:periodic_projective_rigidity}), the $Y_g$
    family assembles into a projective representation of $G$ on the
    bond space: there exist a scalar $2$-cochain
    $\omega : G \times G \to \C^{\times}$ and virtual gauges
    $X : G \to \GL_D(\C)$ satisfying
    \[
        X(g)\, X(h) \;=\; \omega(g, h)\, X(g h),
    \]
    while for each $g$ the matrix $X(g^{-1})$ still realises the
    Corollary~4.1 $\Z_{m_g}$-intertwining between $A$ and the twisted
    tensor $\widetilde{A}_g$.
\end{theorem}

\begin{proof}\leanok
    \uses{cor:cor_4_1_physical_symmetry_zgauge,
          def:periodic_projective_rigidity}
    Unpack the rigidity hypothesis to extract the family
    $(Y, u)$ together with the Corollary-4.1 intertwiners.  Set
    $X(g) := Y(g^{-1})$ and $\omega(g, h) := u(h^{-1}, g^{-1})$.  The
    projective multiplication law $X(g)\, X(h) = \omega(g, h)\, X(gh)$
    is then exactly the $(h^{-1}, g^{-1})$-instance of the factor-system
    identity, using $(gh)^{-1} = h^{-1} g^{-1}$.  The intertwining
    relation between $A$ and $\widetilde{A}_g$ transports unchanged.
\end{proof}

\begin{theorem}[Cocycle identity for the upgraded factor system]%
    \label{thm:cor_4_1_projective_rep_cocycle}
    \lean{MPSTensor.cor_4_1_projective_rep_cocycle}
    \leanok
    \uses{thm:cor_4_1_projective_rep}
    If the bond dimension $D$ is positive, the factor system $\omega$
    produced by Theorem~\ref{thm:cor_4_1_projective_rep} satisfies the
    multiplicative $2$-cocycle identity
    $\omega(g, h)\, \omega(g h, k) = \omega(g, h k)\, \omega(h, k)$.
\end{theorem}

\begin{proof}\leanok
    \uses{thm:cor_4_1_projective_rep}
    Associativity of matrix multiplication applied to
    $X(g)\, X(h)\, X(k)$, combined with the projective multiplication
    law, gives
    \[
        \omega(g, h)\, \omega(g h, k)\, X((g h) k)
        =
        \omega(g, h k)\, \omega(h, k)\, X(g (h k)).
    \]
    Since $(g h) k = g (h k)$, both sides are scalar multiples of the
    same matrix $X((g h) k)$.  The matrix lies in $\GL_D(\C)$ and is
    therefore invertible, but the scalar cancellation leading from an
    equality of scalar multiples of a matrix to equality of the scalar
    coefficients additionally requires $D > 0$ so that the underlying
    index set is inhabited.  Cancelling then yields
    $\omega(g, h)\, \omega(g h, k) = \omega(g, h k)\, \omega(h, k)$
    in $\C^{\times}$, hence in $\mathrm{Units}(\C)$.
\end{proof}

\begin{remark}[SPT classification]%
    \label{rem:cor_4_1_projective_rep_spt}
    \uses{thm:cor_4_1_projective_rep_cocycle}
    The cohomology class of $\omega$ in $H^{2}(G, U(1))$ classifies the
    SPT phase of the symmetric MPS (cf.~Chen--Gu--Wen and
    P\'erez-Garc\'ia et al.).  This classification statement is not
    formalised at present; only the underlying 2-cocycle identity of
    Theorem~\ref{thm:cor_4_1_projective_rep_cocycle} is available in
    the Lean development.
\end{remark}

\begin{theorem}[Bundled projective representation from symmetry]%
    \label{thm:projectiveRep_of_symmetry}
    \lean{MPSTensor.projectiveRep_of_symmetry}
    \leanok
    \uses{thm:cor_4_1_projective_rep, thm:cor_4_1_projective_rep_cocycle}
    Theorems~\ref{thm:cor_4_1_projective_rep}
    and~\ref{thm:cor_4_1_projective_rep_cocycle} are packaged into a
    single high-level statement: from the periodic projective-rigidity
    hypothesis (Definition~\ref{def:periodic_projective_rigidity}), one
    obtains a projective representation of $G$ on the bond space whose
    factor system is a genuine $2$-cocycle whenever $D > 0$.
\end{theorem}

\section{Refinement and divisibility}\label{sec:periodic_refinement}

The next theorem characterises when the matrix-product-vector family of a
periodic tensor admits a finer description in terms of a smaller-period
tensor whose blocks have been merged.  This was the original motivation
for the irreducible-form framework in \cite{DeLasCuevas2017Irreducible}.

\begin{definition}[$p$-divisible CP map]
    \label{def:p_divisible_channel}
    \lean{MPSTensor.IsPDivisibleChannel}
    \leanok
    A linear endomorphism $\mathcal{E}$ of $M_D(\C)$ is \emph{$p$-divisible}
    if there exists a CPTP map $\mathcal{E}'$ of $M_D(\C)$ such that
    $\mathcal{E} = (\mathcal{E}')^p$.
\end{definition}

\begin{definition}[$p$-refinable MPS tensor]
    \label{def:p_refinable}
    \lean{MPSTensor.IsPRefinable}
    \leanok
    \uses{def:mps_tensor, def:blocked_tensor}
    An MPS tensor $B : \mathrm{MPSTensor}\ d\ D$ is \emph{$p$-refinable}
    if there exist another tensor $A : \mathrm{MPSTensor}\ d\ D$ and an
    isometry $W : \C^d \to (\C^d)^{\otimes p}$ (encoded as a matrix
    $W \in M_{d^p \times d}(\C)$ with $W^{\dagger} W = \Id$) such that the
    $p$-blocked tensor $A^{[p]}$ matches the $W$-image of $B$ at the level
    of MPV coefficients:
    \[
        \mathrm{coeff}\big(A^{[p]}\big)\big(\mathrm{ofFn}\,\tau\big)
        \;=\;
        \sum_{\sigma : \mathrm{Fin}\, N \to \mathrm{Fin}\, d}
            \Big(\prod_{k} W_{\tau(k),\,\sigma(k)}\Big)\,
            \mathrm{coeff}\,B\,\big(\mathrm{ofFn}\,\sigma\big)
    \]
    for every $N \in \N$ and every
    $\tau : \mathrm{Fin}\, N \to \mathrm{Fin}\, d^{p}$.

    In paper notation, this encodes
    $\ket{V_{pN}(A)} = W^{\otimes N}\, \ket{V_N(B)}$ for all~$N$.
\end{definition}

\begin{theorem}[Pullback of a refinement witness]
    \label{thm:p_refinement_pullback_irreducible_form}
    \lean{MPSTensor.pRefinementCanonicalization_pullback_of_irreducibleForm}
    \leanok
    \uses{def:irreducible_form, def:p_refinable}
    Let $B$ be an MPS tensor in irreducible form~II. If $B$ is
    $p$-refinable, then there exist a tensor $A$, an isometry
    $W : \C^d \to (\C^d)^{\otimes p}$, and the pullback tensor
    $$
        C^{\alpha} := \sum_{i=0}^{d-1} W_{\alpha,i} B^i
        \qquad (\alpha \in \{0,\ldots,d^p-1\})
    $$
    such that $C$ is again in irreducible form~II, $\mathcal E_C = \mathcal E_B$,
    and $C$ has the same MPV family as $A^{[p]}$.
\end{theorem}

\begin{proof}\leanok
    Choose $(A,W)$ from the refinement data. The pullback construction gives
    $\mathcal E_C = \mathcal E_B$ and equality of MPV families between $C$
    and $A^{[p]}$. To see that $C$ is still in irreducible form~II, transport
    each periodic block $B_j$ of $B$ through the same physical isometry $W$.
    The left-canonical normalization is preserved because $W^{\dagger}W = \Id$,
    while irreducibility and the peripheral spectrum are preserved because the
    transfer map of each transported block agrees with that of $B_j$.
    Reassembling these transported blocks with the same weights gives an
    irreducible-form decomposition of $C$.
\end{proof}

\begin{theorem}[Thm.~4.1 — $p$-refinability $\iff$ $p$-divisibility]%
    \label{thm:thm_4_1_p_refinement}
    \lean{MPSTensor.thm_4_1_p_refinement}
    \leanok
    \uses{def:irreducible_form, def:p_refinable, def:p_divisible_channel,
          thm:periodic_ft, thm:thm_4_1_p_refinement_forward,
          thm:thm_4_1_p_refinement_reverse}
    Let $B$ be an MPS tensor in irreducible form~II and let $p \ge 1$.
    Assume both canonicalization hypotheses used in the one-way statements:
    one for refinability $\Rightarrow$ divisibility and one for divisibility
    $\Rightarrow$ refinability. Then $B$ is $p$-refinable iff its transfer
    map $\mathcal{E}_B$ is $p$-divisible.
\end{theorem}

\begin{proof}\leanok
    \uses{thm:thm_4_1_p_refinement_forward, thm:thm_4_1_p_refinement_reverse}
    Bundle the forward implication
    (Theorem~\ref{thm:thm_4_1_p_refinement_forward}) with the reverse
    implication (Theorem~\ref{thm:thm_4_1_p_refinement_reverse}) into a
    single biconditional.  Both directions are formalised conditionally
    on their respective canonicalization hypotheses, whose combination
    yields the full equivalence.
\end{proof}

\begin{theorem}[Thm.~4.1 forward direction, conditional form]%
    \label{thm:thm_4_1_p_refinement_forward}
    \lean{MPSTensor.thm_4_1_p_refinement_forward}
    \leanok
    \uses{def:irreducible_form, def:p_refinable, def:p_divisible_channel}
    Let $B$ be an MPS tensor in irreducible form~II. Assume a
    \emph{canonicalization} hypothesis providing the remaining analytic
    assumptions (left-canonical witness with matching transfer map). Then
    $p$-refinability of $B$ implies $p$-divisibility of $\mathcal{E}_B$.
\end{theorem}

\begin{proof}\leanok
    The canonicalization hypothesis supplies a witness $A$ with
    $\sum_i (A^i)^\dagger A^i = \Id$ (left-canonical) and
    $\mathcal{E}_B = \mathcal{E}_{A^{[p]}}$. The left-canonical property
    makes $\mathcal{E}_A$ a CPTP channel, and the blocking identity
    $\mathcal{E}_{A^{[p]}} = (\mathcal{E}_A)^p$ (formalised as
    \lean{MPSTensor.transferMap_blockTensor}) then identifies
    $\mathcal{E}_B$ with the $p$-th power of the channel $\mathcal{E}_A$.
\end{proof}

\begin{theorem}[Thm.~4.1 reverse direction, conditional form]%
    \label{thm:thm_4_1_p_refinement_reverse}
    \lean{MPSTensor.thm_4_1_p_refinement_reverse}
    \leanok
    \uses{def:irreducible_form, def:p_refinable, def:p_divisible_channel}
    Let $B$ be an MPS tensor in irreducible form~II and let $p \ge 1$.
    Assume an \emph{inverse canonicalization} hypothesis providing a
    witness $A : \mathrm{MPSTensor}\ d\ D$ with
    $\mathcal{E}_B = \mathcal{E}_{A^{[p]}}$ whenever $\mathcal{E}_B$ is
    $p$-divisible.  Then $p$-divisibility of $\mathcal{E}_B$ implies
    $p$-refinability of $B$.
\end{theorem}

\begin{proof}\leanok
    From the inverse canonicalization hypothesis, extract
    $A$ with $\mathcal{E}_B = \mathcal{E}_{A^{[p]}}$.  This matches two
    finite Kraus families of the same completely positive map: the
    blocked family $A^{[p]}$ with $d^{p}$ operators and the original
    family $B$ with $d$ operators.  The cardinality comparison
    $d \le d^{p}$ holds whenever $p \ge 1$, so Wolf
    Theorem~2.18 (isometric freedom of Kraus representations,
    formalised as \lean{kraus_isometry_freedom_iff}) supplies an
    isometry $V \in M_{d^p \times d}(\C)$ with $V^{\dagger} V = \Id$ and
    $(A^{[p]})^{\alpha} = \sum_{j} V_{\alpha,j}\, B^{j}$.  Expanding
    $\mathrm{coeff}\,(A^{[p]})\,(\mathrm{ofFn}\,\tau)$ by linearity of
    matrix multiplication and of the trace, the contributions
    reorganise into the $V$-weighted sum over length-$N$ words of $B$
    required by Definition~\ref{def:p_refinable}, exhibiting $V$ as the
    isometry $W$ in the statement of $p$-refinability.
\end{proof}

\begin{remark}[Continuum-limit application]\notready
    \label{rem:thm_4_1_continuum_limit}
    The original motivation for Theorem~\ref{thm:thm_4_1_p_refinement}
    is a continuum-limit construction in the spirit of
    \cite{DeLasCuevas2017Irreducible}: refinability along arbitrarily
    large $p$ is equivalent to infinite divisibility of the transfer
    map, which selects out continuum limits of translationally invariant
    MPS.  The further consequences are explored elsewhere.
\end{remark}
